\exercisegroup

\begin{enumerate}
\def\labelenumi{\arabic{enumi})}
\item
  设 E、F 与 G 是三个局部凸空间，S 是 E 的由有界集组成的一个覆盖。证明：若 u 是从 \(E \times F\) 到 G 中的分别连续双线性映射，使得对每个集 \(M \in S\) ，u 在 \(M \times F\) 上的限制连续，则 u 是 S-亚连续的。
\item
  设 \(\mathbf{E}\) 是直和空间 \(\mathbf{R}^{(\mathbf{N})}\) ，赋以 \(\mathbf{R}^{\mathbf{N}}\) 上乘积拓扑所诱导的拓扑。证明：双线性型 \(((x_n), (y_n)) \mapsto \sum_{n=0}^{\infty} x_n y_n\) （在 \(\mathbf{E} \times \mathbf{E}\) 上）是分别连续的，但对每个集 \(\mathfrak{S}\) （由 \(\mathbf{E}\) 的有界子集组成，且至少含有一个无穷维有界集），这个双线性型都不是 \(\mathfrak{S}\) -亚连续的。
\end{enumerate}


\begin{enumerate}
\def\labelenumi{\arabic{enumi})}
\setcounter{enumi}{2}
\tightlist
\item
  设 E 是赋以最细局部凸拓扑（II, 第 26 页）的空间 \(\mathbf{R}^{(\mathbf{N})}\) ；设 F 是空间 \(R^{N}\) ；空间 E 是超有界型的（III, 第 45 页, 习题 19）且完备，而 F 可度量化且完备。设 S（相应地，T）是 E（相应地，F）的所有有界子集组成的集。证明：E × F 上的双线性型 \(((x_{n}), (y_{n})) \mapsto \sum_{n=0}^{\infty} x_{n} y_{n}\) 是 \((\mathfrak{S}, \mathfrak{T})\) -亚连续的，但不连续
\end{enumerate}

（cf.~IV, 第 48 页, 习题 11）。

\begin{enumerate}
\def\labelenumi{\arabic{enumi})}
\setcounter{enumi}{3}
\item
  设 E 是一个局部凸空间，F 是一个次桶型空间（III, 第 44 页, 习题 7），T 是 F 的所有有界子集组成的集。证明：若从 \(E \times F\) 到一个局部凸空间 G 中的双线性映射是 T-亚连续的，则对 E 的有界子集的每个集 S，它都是 \((\mathfrak{S}, \mathfrak{T})\) -亚连续的（cf.~III, 第 44 页, 习题 11）。
\item
  \begin{enumerate}
  \def\labelenumii{\alph{enumii})}
  \tightlist
  \item
    设 \(\mathbf{E}\) 、\(\mathbf{F}\) 与 \(\mathbf{G}\) 是三个 Hausdorff 局部凸空间，\(u\) 是从 \(\mathrm{E} \times \mathrm{F}\) 到 \(\mathbf{G}\) 中的一个双线性映射。要使存在 0 的一个均衡邻域 \(\mathbf{U}\) （\(\mathbf{E}\) 中），使得映射 \(u(x, .)\) （当 \(x\) 取遍 \(\mathbf{U}\) 时）全体所成的集在 \(\mathcal{L}(\mathbf{F}; \mathbf{G})\) 中等度连续，必须且只需：\(u\) 当把 \(\mathbf{E}\) 的拓扑换成使诸集 \(\lambda \mathbf{U} (\lambda \neq 0)\) 构成 0 的基本邻域系的最粗拓扑时是连续的。证明：若 \(\mathbf{G}\) 赋范，则从 \(\mathrm{E} \times \mathrm{F}\) 到 \(\mathbf{G}\) 中的每个连续双线性映射都满足这一条件。
  \end{enumerate}
\end{enumerate}

\begin{enumerate}
\def\labelenumi{\alph{enumi})}
\setcounter{enumi}{1}
\tightlist
\item
  取乘积空间 \(R^{N}\) 作为 E、F 与 G，并取连续双线性映射 \(((x_{n}), (y_{n})) \mapsto (x_{n}y_{n})\) 作为 u。证明：不存在 E 中 0 的任何邻域 U，使得当 x 取遍 U 时的映射 \(u(x, .)\) 全体所成的集在 \(\mathcal{L}(F; G)\) 中等度连续。
\end{enumerate}

\begin{enumerate}
\def\labelenumi{\arabic{enumi})}
\setcounter{enumi}{5}
\tightlist
\item
  设 E、F 与 G 是三个拓扑向量空间。称从 \(E \times F\) 到 G 中的双线性映射的一个集 H 是分别等度连续的，如果对一切 \(x \in E\) ，当 u 取遍 H 时的线性映射 \(u(x, .)\) 全体所成的集在 \(\mathcal{L}(F; G)\) 中等度连续，且对一切 \(y \in F\) ，当 u 取遍 H 时的线性映射 \(u(. , y)\) 全体所成的集在 \(\mathcal{L}(E; G)\) 中等度连续。
\end{enumerate}

假设 F 可度量化，且 E 是 Baire 空间（cf.~III, 第 43 页, 习题 5 与 V, 第 79 页, 习题 15）。证明：从 \(E \times F\) 到 G 中的双线性映射的每个分别等度连续的集都是等度连续的（cf.~III, 第 42 页, 习题 11）。

\begin{enumerate}
\def\labelenumi{\arabic{enumi})}
\setcounter{enumi}{6}
\tightlist
\item
  设 \(\mathrm{E}\) 、\(\mathrm{F}\) 与 \(\mathrm{G}\) 是三个拓扑向量空间，\(\mathfrak{S}\) 是 \(\mathrm{E}\) 的有界子集的一个集，\(\mathrm{H}\) 是从 \(\mathrm{E} \times \mathrm{F}\) 到 \(\mathrm{G}\) 中的分别连续双线性映射的一个集。下列性质等价：
\end{enumerate}

\(\alpha)\) 对 0 的每个邻域 \(\mathbf{W}\) （\(\mathbf{G}\) 中）以及每个集 \(\mathbf{M} \in \mathfrak{S}\) ，存在 0 的一个邻域 \(\mathbf{V}\) （\(\mathbf{F}\) 中），使得 \(u(\mathbf{M} \times \mathbf{V}) \subset \mathbf{W}\) （对一切 \(u \in \mathbf{H}\)）。

\(\beta)\) 对每个集 \(\mathbf{M} \in \mathfrak{S}\) ，\(\mathbf{H} \times \mathbf{M}\) 在映射 \((u, x) \mapsto u(x, .)\) 下的像是 \(\mathcal{L}(\mathrm{F}; \mathrm{G})\) 的一个等度连续子集。

\(\gamma)\) 当 \(u\) 取遍 H 时，映射 \(y \mapsto u(., y)\) （从 F 到 \(\mathcal{L}_{\mathfrak{S}}(\mathrm{E}; \mathrm{G})\) 中）全体所成的集是等度连续的。

这时称 H 是从 \(E \times F\) 到 G 中的（分别连续的）双线性映射的一个 S-等度亚连续集。类似地，对 F 的有界子集的一个集 T，我们定义 T-等度亚连续集与 \((\mathfrak{S}, \mathfrak{T})\) -等度亚连续集的概念。

\begin{enumerate}
\def\labelenumi{\arabic{enumi})}
\setcounter{enumi}{7}
\item
  设 H 是从 \(E \times F\) 到 G 中的双线性映射的一个 S-等度亚连续集（习题 7）。对每个子集 \(M \in S\) ，证明 H 在 \(M \times F\) 上等度连续；此外，对 F 的每个有界子集 Q，诸集 \(u(\mathbf{M} \times \mathbf{Q})\) （其中 u 取遍 H）的并在 G 中有界。
\item
  设 H 是双线性映射的一个 \((\mathfrak{S}, \mathfrak{T})\) -等度亚连续集（从 \(E \times F\) 到 G 中）（习题 7）；证明：对每一对集 \(M \in S\) 、\(N \in T\) ，H 在 \(M \times N\) 上一致等度连续。
\item
  设 \(\mathrm{E}_1\) 、\(\mathrm{E}_2\) 与 \(\mathrm{F}\) 是三个拓扑向量空间，\(\mathbf{G}_1\) （相应地，\(\mathbf{G}_2\) ）是 \(\mathrm{E}_1\) （相应地，\(\mathrm{E}_2\) ）的一个处处稠密子空间，\(\mathfrak{S}_1\) （相应地，\(\mathfrak{S}_2\) ）是 \(\mathbf{G}_1\) （相应地，\(\mathbf{G}_2\) ）的有界子集的一个族。设 \(\mathrm{H}\) 是从 \(\mathrm{E}_1 \times \mathrm{E}_2\) 到 \(\mathrm{F}\) 中的分别连续双线性映射的一个集；若在 \(\mathbf{G}_1 \times \mathbf{G}_2\) 上取诸映射 \(u \in \mathrm{H}\) 的限制，所得之集是 \((\mathfrak{S}_1, \mathfrak{S}_2)\) -等度亚连续的，则 \(\mathrm{H}\) 亦是。
\item
  若 \(\mathbf{F}\) 是桶型空间，则从 \(\mathbf{E} \times \mathbf{F}\) 到一个局部凸空间 \(\mathbf{G}\) 中的双线性映射的每个分别等度连续的集都是 \(\mathfrak{S}\) -等度亚连续的（对每个集 \(\mathfrak{S}\) ，由 \(\mathbf{E}\) 的有界子集组成）。
\item
  设 E、F 是两个拓扑向量空间，f 是双线性映射 \((x, u) \mapsto u(x)\) （从 \(E \times \mathcal{L}(E; F)\) 到 F 中）；设 T 是 \(\mathcal{L}(E; F)\) 上与向量空间结构相容且比简单收敛拓扑更细的一个拓扑。设 S 是 E 的有界子集的一个族，U 是 \(\mathcal{L}(E; F)\) 的（对拓扑 T 而言的）有界子集的一个族。证明：f 是 S-亚连续的，当且仅当 T 比 S-拓扑细；f 是 U-亚连续的，当且仅当 U 中诸集是 \(\mathcal{L}(E; F)\) 的等度连续子集。
\item
  设 E、F、G 是三个拓扑向量空间，S（相应地，T）是 E（相应地，F）的有界子集的一个族。设 H 是从 \(E \times F\) 到 G 中的 T-亚连续双线性映射所组成的向量空间。
\end{enumerate}

\begin{enumerate}
\def\labelenumi{\alph{enumi})}
\item
  证明：在 H 上，对形如 \(M \times N\) （其中 \(M \in S\) 、\(N \in T\) ）的集合的一致收敛拓扑与向量空间结构相容；这个拓扑称为 H 上的 \((\mathfrak{S}, \mathfrak{T})\) -拓扑。对每个映射 \(u \in H\) ，令 \(\tilde{u}\) 为连续映射 \(x \mapsto u(x, .)\) （从 E 到 \(\mathcal{L}_{\mathfrak{T}}(F; G)\) 中）。证明 \(u \mapsto \tilde{u}\) 是从赋以 \((\mathfrak{S}, \mathfrak{T})\) -拓扑的空间 H 到空间 \(\mathcal{L}_{\mathfrak{S}}(E; \mathcal{L}_{\mathfrak{T}}(F; G))\) 上的一个同构。
\item
  设 L 是 H 的一个子集，使得对每一对 \((x, y) \in E \times F\) ，当 u 取遍 L 时的 \(u(x, y)\) 全体所成的集在 G 中有界（H 的简单有界子集）。证明：若 E、F、G 是局部凸的，且 E 与 F 是 Hausdorff 的和拟完备的，则 L 在 H 中对 \((\mathfrak{S}, \mathfrak{T})\) -拓扑有界。
\item
  设 E、F、G 是三个 Hausdorff 局部凸空间。若 E 是桶型的且 F 拟完备或桶型，则 H 的每个简单有界子集 L 都是 T-等度亚连续的（III, 第 47 页, 习题 7）。
\end{enumerate}

\(d\) ) 若 E 与 F 是桶型的，G 拟完备，且 \(\mathfrak{S}\) 与 \(\mathfrak{T}\) 分别是 E 与 F 的覆盖，则 H 对 \((\mathfrak{S},\mathfrak{T})\) -拓扑是 Hausdorff 的且拟完备的。

\begin{enumerate}
\def\labelenumi{\arabic{enumi})}
\setcounter{enumi}{13}
\item
  把 § 5 的定义与结果推广到任意多重线性映射。设 E、F、G 是三个拓扑向量空间，S（相应地，T）是 E（相应地，F）的有界子集的一个族，U 是从 E × F 到 G 中的双线性 (S, T)-亚连续映射所组成的空间 \(\mathcal{L}_{\mathfrak{S},\mathfrak{T}}\) (E, F; G)（赋以 (S, T)-拓扑（习题 13））的有界子集的一个族。证明：从 E × F × \(\mathcal{L}_{\mathfrak{S},\mathfrak{T}}\) (E, F; G) 到 G 中的三线性映射 (x, y, u) → u(x, y) 是 (S, T)-亚连续的；要使它是 (S, U)-亚连续的，必须且只需 U 中的每个集 L 都是 S-等度亚连续的（III, 第 47 页, 习题 7）。
\item
  设 \(E\) 是由所有实数序列 \(x = (\xi_n)_{n \geqslant 0}\) （使通项为 \(\xi_n\) 的级数收敛）组成的空间。令 \(\| x \| = \sup_n | \sum_{k=0}^n \xi_k |\) 。
\end{enumerate}

\begin{enumerate}
\def\labelenumi{\alph{enumi})}
\tightlist
\item
  证明 \(\| x\|\) 是 E 上的一个范数，且 E 对这个范数是完备的。
\end{enumerate}

\begin{enumerate}
\def\labelenumi{\alph{enumi})}
\setcounter{enumi}{1}
\tightlist
\item
  证明：向量空间 \(\ell^1 (\mathbf{N})\) （I, 第 4 页）作为 E 的子空间（对 E 的拓扑而言）是处处稠密的；\(\ell^1 (\mathbf{N})\) 上由范数 \(\| x\| _1 = \sum_{n = 0}^{\infty}|\xi_n|\) 定义的拓扑严格地比 E 的拓扑所诱导的拓扑细。
\end{enumerate}

\begin{enumerate}
\def\labelenumi{\alph{enumi})}
\setcounter{enumi}{2}
\tightlist
\item
  设 \((\mathbf{P}_n)\) 是 \(\mathbf{N} \times \mathbf{N}\) 的有限子集的递增序列，构成 \(\mathbf{N} \times \mathbf{N}\) 的一个覆盖。对每个 \(x = (\xi_n) \in \mathrm{E}\) 及每个 \(y = (\eta_n) \in \ell^1(\mathbf{N})\) ，令 \(f_n(x, y) = \sum_{(i,j) \in \mathbf{P}_n} \xi_i \eta_j\) 。那么，序列 \((f_n(x, y))\) 对每对 \((x, y) \in \mathrm{E} \times \ell^1(\mathbf{N})\) 都趋于一个极限，当且仅当对这些对 \((x, y)\) 中的每一对，序列 \((f_n(x, y))\) 有界；此时 \(f_n(x, y)\) 的极限等于 \(\left(\sum_{n=0}^{\infty} \xi_n\right) \left(\sum_{n=0}^{\infty} \eta_n\right)\) 。
\end{enumerate}

（利用 III, 第 48 页, 习题 13, c)，证明双线性型的序列 \((f_n)\) 等度连续，并注意它在子空间 \(\ell^1 (\mathbf{N})\times \ell^1 (\mathbf{N})\) 中收敛；再利用 \(b).\) 完成论证。）\(d)\) 对每个 \(j\in \mathbf{N}\) ，令 \(\rho_{jn}\) 为 \(\mathbf{N}\) 的闭区间的最小个数，使这些闭区间的并是 \(\mathbf{P}_n\cap (\mathbf{N}^{\prime}\times \{j\})\) 在 \(\mathbf{N}\) 上的投影；令 \(\rho_{n} = \sup_{j\in \mathbf{N}}\rho_{jn}\) 。证明在 \(c)\) 中得到的条件等价于 \(\sup_n\rho_n < + \infty\) 。（若 \(\phi_n\) 是 \(\mathbf{P}_n\) 的特征函数，证明

双线性型 \(f_{n}\) 的范数是 \(\sup_{j\in \mathbf{N}}\big(\sum_{i = 0}^{\infty}|\rho_n(i,j) - \phi_n(i + 1,j)|\big).\)

